% Preambulo comun de documentacion.tex y manual_usuario.tex
\documentclass[11pt,a4paper]{article}

\usepackage{fontspec}
\usepackage{polyglossia}
\setmainlanguage{spanish}
\setmainfont{Liberation Serif}
\setmonofont{Liberation Mono}[Scale=0.9]
\usepackage[margin=2.5cm]{geometry}
\usepackage[table]{xcolor}
\usepackage{graphicx,fancyhdr,longtable,array,listings,caption,float}
\usepackage[hidelinks]{hyperref}
\Urlmuskip=0mu plus 2mu\def\UrlBreaks{\do\/\do-\do.}

\setlength{\parskip}{0.7em}
\setlength{\parindent}{0pt}
\setlength{\emergencystretch}{2em}
\linespread{1.15}
\renewcommand{\arraystretch}{1.2}

% "1. Introduccion" con punto solo en las secciones; "3.1 Titulo" sin punto
\makeatletter
\def\@seccntformat#1{\csname the#1\endcsname\csname punto@#1\endcsname\quad}
\def\punto@section{.}
\makeatother
\addto\captionsspanish{\renewcommand{\contentsname}{Tabla de contenidos}}

% "Pagina X de Y" sin el paquete lastpage
\AtEndDocument{\label{ultimapagina}}
\pagestyle{fancy}
\fancyhf{}
\renewcommand{\headrulewidth}{0pt}
\fancyfoot[C]{\small Página \thepage\ de \pageref{ultimapagina}}

\newcommand{\kw}[1]{\texttt{#1}}
\newcommand{\ggez}{GGEZ}

\lstdefinelanguage{GGEZ}{
  sensitive=true,
  morecomment=[l]{/w},
  morecomment=[s]{/all}{/end},
  morestring=[b]",
}
\lstset{
  language=GGEZ, basicstyle=\ttfamily\small, columns=fullflexible, keepspaces=true,
  breaklines=true, showstringspaces=false, xleftmargin=1.2em, aboveskip=0.4em, belowskip=0.6em,
}

% Tablas al estilo del documento: encabezado gris y filas de categoria
\newcommand{\cab}{\rowcolor{gray!22}}
\newcommand{\categoria}[2]{\multicolumn{#1}{|l|}{\cellcolor{gray!10}\textbf{#2}}\\\hline}

% Tabla de pruebas: la generan Traductor/pruebas/correr.py
\newenvironment{tablaprueba}{\footnotesize
  \begin{longtable}{|p{3.4cm}|p{3.0cm}|p{5.6cm}|p{2.1cm}|}\hline
  \cab\textbf{Prueba y objetivo} & \textbf{Entrada} & \textbf{Resultado esperado} & \textbf{Resultado obtenido y estado}\\\hline\endhead}
  {\end{longtable}}
